\chapter{Vaginal Bleeding in Pregnancy}

\section*{Definition}
Vaginal bleeding during pregnancy ranges from light spotting to heavy hemorrhage. Some early spotting is benign, but bleeding can indicate miscarriage, ectopic pregnancy, placental problems, or labor.

\section*{When to See a Doctor}
Emergency care is required for heavy bleeding, severe abdominal pain, fainting, shoulder pain, dizziness, fever, contractions, fluid leakage, or reduced fetal movement. Arrange prompt evaluation for persistent, recurrent, unexplained, or function-limiting symptoms.

\section*{How It Develops}
The symptom results from irritation, inflammation, altered organ function, injury, infection, or disruption of normal neurologic or vascular signaling. Age, pregnancy status, onset, distribution, severity, and associated findings determine the urgency and likely causes.

\section*{Causes}
\begin{itemize}
  \item Implantation or cervical bleeding, threatened miscarriage, ectopic pregnancy, placenta previa, placental abruption, infection, and labor-related cervical change.
  \item Medication effects, environmental exposures, and underlying chronic disease may also contribute.
  \item Serious causes are less common but must be considered when symptoms are sudden, progressive, severe, or associated with systemic illness.
\end{itemize}

\section*{Related Symptoms}
\begin{itemize}
  \item Pain, fever, fatigue, nausea, weakness, or changes in normal function
  \item Bleeding, swelling, discharge, breathing difficulty, or neurologic symptoms when a serious cause is present
\end{itemize}

\section*{Medical Evaluation}
\begin{itemize}
  \item History addresses onset, duration, severity, triggers, injuries, exposures, medications, medical conditions, age, and pregnancy status when relevant.
  \item Examination includes vital signs and focused assessment of the relevant organ system, neurologic status, hydration, and function.
  \item Testing may include blood or urine studies, cultures, imaging, physiologic testing, fetal or pediatric assessment, or specialist examination.
\end{itemize}

\section*{Treatment Options}
Treatment is directed at the cause and may include supportive care, targeted medication, treatment of infection or inflammation, correction of metabolic disease, physical therapy, or specialist procedures. Persistent symptoms require reassessment.

\section*{Self-Care and Home Management}
Avoid known triggers, protect the affected area, maintain appropriate hydration, and record timing and associated symptoms. Use only age- or pregnancy-appropriate medicines recommended by a clinician. Do not delay emergency care while trying home treatment.

\section*{References}
\begin{thebibliography}{99}
\bibitem{vaginal_bleeding_in_pregnancy:acog} American College of Obstetricians and Gynecologists. Vaginal Bleeding in Early Pregnancy. ACOG Practice Bulletin No. 176. \textit{Obstet Gynecol}. 2017;129:e78--e89.
\bibitem{vaginal_bleeding_in_pregnancy:rcog-ectopic} Royal College of Obstetricians and Gynaecologists. Diagnosis and Management of Ectopic Pregnancy. RCOG Green-top Guideline No. 21; 2016.
\bibitem{vaginal_bleeding_in_pregnancy:rcog-aph} Royal College of Obstetricians and Gynaecologists. Antepartum Haemorrhage. RCOG Green-top Guideline No. 63; 2011 (updated 2021).
\bibitem{vaginal_bleeding_in_pregnancy:nice} National Institute for Health and Care Excellence. Antenatal care. NICE guideline NG201; 2021 (updated 2023).
\end{thebibliography}
